\chapter{Is AI a Bubble?}

Nowadays there is a debate about whether AI is a bubble.
\section{Alibaba's Rationalization}

The following is from Alibaba's 2025 Q3 earnings call.
\begin{quote}
First is from the demand side. On the demand side, we see two very obvious drivers. One is that we
see base models, whether foundation models we see, or video generation models, or more future full-modal
models, capabilities in these areas are continuously improving. The entire scaling law trend still hasn’t hit a
wall. We see many breakthroughs from Gemini 3’s progress.

As model capabilities improve, AI models can do more and more things, and can adapt to more and
more scenarios. In the process of improving capabilities for adapted scenarios or tasks that can be driven,
this is also accompanied by the penetration rate of these tasks in various industries. So driven by these two
forces, within three years we see AI demand is still very certain.

Under such high growth demand, we look at the supply side. You analysts may have also seen that starting
from the second half of this year, globally, whether from wafer manufacturers, to DRAM manufacturers, to
storage manufacturers, to CPU and other AI server links, all are in short supply. This wave of shortages is
all driven by AI demand. For future expansion cycles of supply chain manufacturers in all aspects, we know
they have considerable bottlenecks. But in this, we think at least it’s a two- to three-year expansion cycle.
Within these two to three years of expansion cycle, we see demand growth, but the supply side is difficult to
rapidly improve. So we think within these three years, AI resources are still in a state of supply falling short
of demand. And we actually see that in our industry, including ourselves, including the U.S. hyperscalers we
see, not only are their new GPUs basically running at full capacity, but even previous generations, or CPUs
from three or five years ago, are running at full capacity. So we think within three years, the so-called AI
bubble should not really exist.
\end{quote}

\section{Buffett's Grocery Store Escalator Story}

Warren Buffett once owned a small grocery store that sat across the street from a rival. One day, the competitor installed an escalator to impress customers. Buffett felt he had no choice but to follow and put in an escalator of his own. In the end, nothing important changed: the same people shopped on each side of the street as before, but now both stores had higher costs and no real gain in business.

Buffett’s lesson is that copying a rival’s flashy investment just to ``keep up'' often leaves everyone worse off if it doesn’t create lasting value for customers—a useful analogy for today’s AI spending race.\nocite{buffett_escalator_ai}

\section{What Is the Value of AI?}

Alibaba rationalizes that AI is a high-growth area simply because everyone else is investing and demand appears to be rising.

He argues that AI is not a bubble because there is a supply shortage. But the same shortage can also be described as a sign of a bubble. In logical form, his argument reduces to: ``AI is not a bubble because it is not a bubble.''

The speaker says, ``Within these two to three years of expansion cycle, we see demand growth.'' This is speculation about the future. He does not mention at all what concrete value AI will create for customers.

Even if his speculation is correct and demand grows in two or three years, new hardware will also arrive in two or three years. So what is the point of burning money to train large models on today's hardware when tomorrow's hardware will be cheaper and more powerful, and the real demand may be clearer?

It is also unclear what exactly the company is investing in. Is it \vocab{language models}{Language Model}{A language model is a probabilistic model designed to predict the likelihood of a sequence of words. It assigns probabilities to sequences of words, helping to determine the most likely next word in a given context.}? Everyone is training their own language models, and Alibaba does not design its own GPUs. So for what, precisely, is Alibaba burning money? 


